\chapter{Αποτελέσματα}

\section{Κεντρική Εκπαίδευση}

Στην κεντρική εκπαίδευση το  μοντέλο εκπαιδεύτηκε με πρόσβαση σε ολόκληρο το σύνολο εκπαίδευσης. Η συνθήκη αυτή χρησιμοποιείται ως σημείο αναφοράς για την αξιολόγηση της απόδοσης του μοντέλου, καθώς δεν εισάγει τους περιορισμούς της ομοσπονδιακής μάθησης, όπως ο κατακερματισμός των δεδομένων σε \en{clients} και η συνάθροιση τοπικών ενημερώσεων. Το μοντέλο που χρησιμοποιήθηκε ήταν το \en{bert-base-uncased} με προσαρμογή \en{LoRA}, ενώ η αξιολόγηση έγινε στο κοινό σύνολο επικύρωσης και, μετά την επιλογή του καλύτερου συνδυασμού υπερπαραμέτρων, στο σύνολο ελέγχου.

Για την επιλογή των υπερπαραμέτρων πραγματοποιήθηκε αναζήτηση πλέγματος (\en{grid search}). Οι παράμετροι που εξετάστηκαν ήταν ο ρυθμός μάθησης, το μέγεθος παρτίδας εκπαίδευσης και ο βαθμός της \en{LoRA} προσαρμογής. Ο συνολικός χώρος αναζήτησης περιλάμβανε \(3 \times 2 \times 3 = 18\) διαφορετικούς συνδυασμούς, όπως φαίνεται στον Πίνακα~\ref{tab:central_grid_search_space}.

\begin{table}[H]
    \centering
    \caption{Χώρος αναζήτησης υπερπαραμέτρων για την κεντρική εκπαίδευση.}
    \label{tab:central_grid_search_space}
    \begin{tabular}{|l|c|}
        \hline
        \textbf{Υπερπαράμετρος} & \textbf{Τιμές που εξετάστηκαν} \\
        \hline
        \en{Learning rate} & \(2\cdot10^{-5}, 5\cdot10^{-5}, 10^{-4}\) \\
        \hline
        \en{Train batch size} & \(32, 64\) \\
        \hline
        \en{LoRA rank} & \(8, 16, 32\) \\
        \hline
    \end{tabular}
\end{table}

Κριτήριο επιλογής αποτέλεσε η επίδοση στο σύνολο επικύρωσης, με έμφαση στο \en{F1-score}, καθώς το πρόβλημα αφορά δυαδική ταξινόμηση ανωμαλιών και η ισορροπία μεταξύ \en{precision} και \en{recall} είναι σημαντικότερη από την απλή ακρίβεια. Ο καλύτερος συνδυασμός που προέκυψε από την αναζήτηση ήταν ρυθμός μάθησης \(5\cdot10^{-5}\), μέγεθος παρτίδας \(64\) και \en{LoRA rank} ίσο με \(16\). Οι υπόλοιπες βασικές ρυθμίσεις του τελικού πειράματος παρουσιάζονται στον Πίνακα~\ref{tab:central_best_params}.

\begin{table}[H]
    \centering
    \caption{Επιλεγμένες υπερπαράμετροι για το τελικό κεντρικό πείραμα.}
    \label{tab:central_best_params}
    \begin{tabular}{|l|c|}
        \hline
        \textbf{Παράμετρος} & \textbf{Τιμή} \\
        \hline
        \en{Model} & \en{bert-base-uncased} \\
        \hline
        \en{Learning rate} & \(5\cdot10^{-5}\) \\
        \hline
        \en{Train batch size} & \(64\) \\
        \hline
        \en{Evaluation batch size} & \(32\) \\
        \hline
        \en{LoRA rank} & \(16\) \\
        \hline
        \en{LoRA alpha} & \(32\) \\
        \hline
        \en{LoRA dropout} & \(0.1\) \\
        \hline
        Εποχές εκπαίδευσης & \(10\) \\
        \hline
        Σταθμισμένη συνάρτηση κόστους & Ναι \\
        \hline
    \end{tabular}
\end{table}

Η πορεία της εκπαίδευσης στο σύνολο επικύρωσης φαίνεται στον Πίνακα~\ref{tab:central_validation_results}. Η επίδοση βελτιώνεται σημαντικά στις πρώτες εποχές και στη συνέχεια σταθεροποιείται. Η καλύτερη τιμή \en{F1-score} παρατηρήθηκε στην 9η εποχή και παρέμεινε ίδια στην 10η, με \en{F1-score} \(0.9502\), \en{precision} \(0.9721\) και \en{recall} \(0.9320\). Η σταθερότητα των τελευταίων εποχών δείχνει ότι το μοντέλο έχει συγκλίνει χωρίς έντονη υποβάθμιση της απόδοσης στο σύνολο επικύρωσης.

\begin{table}[H]
    \centering
    \caption{Αποτελέσματα επικύρωσης ανά εποχή για το τελικό κεντρικό πείραμα.}
    \label{tab:central_validation_results}
    \begin{tabular}{|c|c|c|c|c|c|}
        \hline
        \textbf{Εποχή} & \textbf{\en{Loss}} & \textbf{\en{Accuracy}} & \textbf{\en{Precision}} & \textbf{\en{Recall}} & \textbf{\en{F1}} \\
        \hline
        1 & 0.2790 & 0.9401 & 0.9163 & 0.9105 & 0.9134 \\
        \hline
        2 & 0.2398 & 0.9592 & 0.9564 & 0.9245 & 0.9392 \\
        \hline
        3 & 0.1878 & 0.9630 & 0.9628 & 0.9296 & 0.9449 \\
        \hline
        4 & 0.1984 & 0.9641 & 0.9644 & 0.9311 & 0.9465 \\
        \hline
        5 & 0.2201 & 0.9614 & 0.9665 & 0.9216 & 0.9417 \\
        \hline
        6 & 0.2009 & 0.9630 & 0.9686 & 0.9244 & 0.9442 \\
        \hline
        7 & 0.2352 & 0.9603 & 0.9678 & 0.9174 & 0.9397 \\
        \hline
        8 & 0.1748 & 0.9662 & 0.9707 & 0.9317 & 0.9494 \\
        \hline
        9 & 0.1821 & 0.9668 & 0.9721 & 0.9320 & 0.9502 \\
        \hline
        10 & 0.1856 & 0.9668 & 0.9721 & 0.9320 & 0.9502 \\
        \hline
    \end{tabular}
\end{table}

Μετά την επιλογή των υπερπαραμέτρων, το τελικό μοντέλο αξιολογήθηκε στο σύνολο ελέγχου. Τα αποτελέσματα παρουσιάζονται στον Πίνακα~\ref{tab:central_test_results}. Το μοντέλο πέτυχε \en{accuracy} \(0.9663\) και \en{F1-score} \(0.9494\), τιμές πολύ κοντά στις αντίστοιχες τιμές επικύρωσης. Αυτό δείχνει ότι η απόδοση του μοντέλου γενικεύεται ικανοποιητικά σε μη χρησιμοποιημένα δεδομένα.

\begin{table}[H]
    \centering
    \caption{Τελικές μετρικές στο σύνολο ελέγχου για την κεντρική εκπαίδευση.}
    \label{tab:central_test_results}
    \begin{tabular}{|l|c|}
        \hline
        \textbf{Μετρική} & \textbf{Τιμή} \\
        \hline
        \en{Test loss} & 0.1884 \\
        \hline
        \en{Accuracy} & 0.9663 \\
        \hline
        \en{Precision} & 0.9717 \\
        \hline
        \en{Recall} & 0.9308 \\
        \hline
        \en{F1-score} & 0.9494 \\
        \hline
    \end{tabular}
\end{table}

Συνολικά, η κεντρική εκπαίδευση οδηγεί σε ισχυρή επίδοση για την ανίχνευση ανωμαλιών στο \en{HDFS}. Η υψηλή τιμή \en{precision} δείχνει ότι οι περισσότερες ακολουθίες που χαρακτηρίζονται ως ανώμαλες είναι πράγματι ανωμαλίες, ενώ η υψηλή τιμή \en{recall} δείχνει ότι το μοντέλο εντοπίζει το μεγαλύτερο μέρος των πραγματικών ανωμαλιών. Επομένως, τα αποτελέσματα της κεντρικής εκπαίδευσης αποτελούν κατάλληλο άνω σημείο αναφοράς για τη σύγκριση με τα ομοσπονδιακά πειράματα.

\section{Ομοσπονδιακή Εκπαίδευση}